% models/mathematical/galaktische_potentiale.tex
% Galaktische Potentiale und ihre Rolle im integrierten Zeitfeldmodell

\section{Galaktische Potentiale im Zeitfeldmodell}

\subsection{Grundidee}

Die galaktische Schicht des Zeitfeldes beschreibt großskalige
Zeitgradienten, die durch:

\begin{itemize}
\item Dunkle-Materie-Halos,
\item Spiralarmstrukturen,
\item vertikale und radiale Oszillationen des Sonnensystems,
\item Supernova-Rate und kosmische Strahlung
\end{itemize}

erzeugt werden.

Diese Strukturen beeinflussen die Zeitrate über gravitative und
strahlungsinduzierte Modulationen.

\subsection{Dunkle-Materie-Halo: NFW-Profil}

Das Navarro–Frenk–White-Profil:
\[
\rho(r) = \frac{\rho_0}{\frac{r}{r_s}\left(1 + \frac{r}{r_s}\right)^2}.
\]

Das zugehörige Potential:
\[
\Phi_{\text{NFW}}(r) = -4\pi G \rho_0 r_s^2 \frac{\ln(1 + r/r_s)}{r/r_s}.
\]

Die Zeitfeldkomponente:
\[
\mathcal{T}^{(g)}_{00}(r) = \Phi_{\text{NFW}}(r).
\]

\subsection{Spiralarm-Potential}

Ein vereinfachtes logarithmisches Spiralarm-Potential:
\[
\Phi_{\text{sp}}(R,\phi,z) =
A_{\text{sp}} \cos\left(m\phi - \frac{\ln(R/R_0)}{\tan p}\right)
e^{-|z|/h_z}.
\]

Parameter:
\begin{itemize}
\item $m$ = Anzahl der Spiralarme,
\item $p$ = Pitch-Winkel,
\item $h_z$ = vertikale Skalenhöhe.
\end{itemize}

\subsection{Vertikale Oszillation des Sonnensystems}

Die Bewegung senkrecht zur galaktischen Ebene:
\[
z(t) = Z_0 \cos(\omega_z t + \phi_z).
\]

Das zugehörige Potential:
\[
\Phi_z(t) = \frac{1}{2} \nu^2 z(t)^2.
\]

\subsection{Radiale Oszillation}

Die radiale Bewegung:
\[
R(t) = R_0 + \Delta R \cos(\omega_R t + \phi_R).
\]

Potential:
\[
\Phi_R(t) = \frac{1}{2} \kappa^2 (R(t) - R_0)^2.
\]

\subsection{Supernova-Rate und kosmische Strahlung}

Die kosmische Strahlungsflussdichte:
\[
J_{\text{CR}}(t) = J_0 + \Delta J \cos(\omega_{\text{CR}} t + \phi_{\text{CR}}).
\]

Die Zeitfeldkopplung:
\[
\mathcal{T}^{(g)}_{0i} \propto J_{\text{CR}}(t).
\]

\subsection{Gesamtpotential der galaktischen Schicht}

Das vollständige galaktische Potential:
\[
\Phi_{\text{gal}}(x) =
\Phi_{\text{NFW}}(r)
+ \Phi_{\text{sp}}(R,\phi,z)
+ \Phi_z(t)
+ \Phi_R(t).
\]

\subsection{Beitrag zum Zeitfeldtensor}

Die galaktische Schicht trägt bei:
\[
\mathcal{T}^{(g)}_{00}(x) = \Phi_{\text{gal}}(x),
\]
\[
\mathcal{T}^{(g)}_{ij}(x) = \partial_i \partial_j \Phi_{\text{gal}}(x).
\]

\subsection{Interferenz mit solaren und planetaren Feldern}

Die Kopplung erfolgt über:
\[
\mathcal{K}\big(\mathcal{T}^{(g)}, \mathcal{T}^{(s)}\big)
= \lambda_{gs} \, \Phi_{\text{gal}} \, A_{\odot}(t),
\]
\[
\mathcal{K}\big(\mathcal{T}^{(g)}, \mathcal{T}^{(p)}\big)
= \lambda_{gp} \, \Phi_{\text{gal}} \sum_k A_k \cos(\omega_k t).
\]

\subsection{Ergebnis}

Die galaktische Schicht definiert die großskaligen Zeitgradienten und
moduliert die solaren und planetaren Zeitfeldkomponenten.
